% Official solution, 2023 theory solutions pp. 22-24. The interleaved
% marking-rubric tables ("Noticing that R1 = R2 ... 1 point", etc.) are
% omitted per transcription policy.

For each of the systems, since the eclipses are short we can assume that the
angle swept along the orbit during the eclipse is negligibly small, which means
that the tangential velocity in relative orbital motion is constant.

\medskip
\noindent\textbf{System A}

A sharp peak minimum in a central eclipse implies that $R_1 = R_2$. Therefore,
both eclipses will be total, with minima of brightnesses $m_1$ and $m_2$,
respectively.

Define the magnitude at minimum as $m_{\min} = m_1$ and the baseline magnitude
as $m_{1,2} = m_1 + m_2$.
% sic — the source writes "m_{1,2} = m_1 + m_2"; m_{1,2} is really the
% combined magnitude of both stars, not a literal sum of magnitudes.
From Pogson's law applied to the two cases:
$m_1 - m_{1,2} = -2.5\log_{10}\left(L_1/(L_1 + L_2)\right)$ and
$m_2 - m_1 = -2.5\log_{10}\left(L_2/L_1\right)$. For the first case, with
simple operations we obtain: $L_2/L_1 = 10^{(m_1 - m_{1,2})/2.5} - 1$.
Plugging this into the second case:
\[ m_2 = m_1 - 2.5\log_{10}\left(10^{(m_1 - m_{1,2})/2.5} - 1\right). \]

As the orbits are near-circular, timescales will be the same for both eclipses.
Therefore it is sufficient to `copy' the shape of the primary one while scaling
down the depth to $|m_{1,2} - m_2|$. Reading off $m_{1,2}$ and $m_{\min}$ from
the plot, we arrive at $L_2/L_1 = 2$, $m_2 = 11.44$, and the following figure:

% FIGURE NEEDED: predicted secondary-eclipse light curve for system A (Bolek) —
% same V-shaped profile and contact times as the primary but with depth scaled
% to minimum V = 11.44 mag (2023 theory solutions, Figure 5, page 22).

\medskip
\noindent\textbf{System B}

For the second system, we follow a generally similar scheme, but we see that
the minimum is flat. $t_{\mathrm{fall}}$ is the time of the descent from the
baseline brightness to the lowest point, i.e., between contacts I and II, while
$t_{\mathrm{flat}}$ is the time between contacts II and III during which the
lightcurve remains at its lowest point:

% FIGURE NEEDED: schematic trapezoidal eclipse light curve marking contacts
% I-IV, with t_fall between contacts I and II and t_flat between contacts II
% and III (2023 theory solutions, page 23).

From this, we obtain
$t_{\mathrm{flat}}/t_{\mathrm{fall}} = 2(R_2 - R_1)/2R_1$ and, after
rearranging, $R_2/R_1 = t_{\mathrm{flat}}/t_{\mathrm{fall}} + 1$.

Here $R_1$ is the radius of the smaller star and there are two solutions --- it
could either be (partially) \emph{eclipsing} or (entirely) \emph{eclipsed}.
Let us first assume that it was eclipsed. After measuring $m_{\min}$ for the
eclipse and $m_{1,2}$ for the baseline we again (just like for system A) assume
$m_{\min} = m_1$ and write down: $L_1/L_2 = 10^{(m_1 - m_{1,2})/2.5} - 1$.

Reading off $m_{1,2}$ and $m_{\min}$ from the plot, we arrive at $L_2/L_1 = 4$,
$R_2/R_1 = 2$.\\
$L_i = S_i \pi R_i^2$, where $S_i$ stands for the surface brightness of each
star.\\
$\implies S_1 = S_2 = S$ (or
$T_{\mathrm{eff},1} = T_{\mathrm{eff},2} = T_{\mathrm{eff}}$).

In such a case, \emph{it does not matter} if star 2 is eclipsing or eclipsed.
During the baseline, we always see two disks with a total brightness of
$S(\pi R_1^2 + \pi R_2^2)$, while during the eclipse, we see one disk of radius
$R_2$ and total brightness $S\pi R_2^2$.

\textbf{Note about marking:} This solution can be independently verified by
assuming that the smaller component was eclipsing and cutting out a disk of
surface $\pi R_1^2$ and brightness contribution $S_1 \pi R_1^2$ in the star of
surface $\pi R_2^2$ and surface brightness $S_2$. Instead of
$m_{\min} = m_1$ of star 1, we would have measured
$m_{\min} = m_{\mathrm{ecl},1}$ of a combination of the two, with a brightness
$L_{ecl1} = S_1 \pi R_1^2 + S_2 \pi (R_2^2 - R_1^2)$. The same conclusion
$S_1 = S_2 = S$ will follow. Any solution correctly arriving at this conclusion
and the correct figure --- no matter the order of steps taken --- should be
scored equally. For the maximum number of points, however, the student should
make a convincing point that this is the \emph{only} solution, i.e.\ lifting
the previously made assumption (or exploring both cases).

To draw the second eclipse, one must simply copy the shape of the first one:

% FIGURE NEEDED: predicted secondary-eclipse light curve for system B (Lolek) —
% identical trapezoidal shape, depth and contact times as the primary eclipse
% (2023 theory solutions, Figure 6, page 24).
